\chapter{PagedAttention 与 vLLM}
\label{chap:paged-attention-vllm}

\section{Serving 系统中的内存碎片问题}
\label{sec:serving-memory-fragmentation}

上一章我们讨论了 KV Cache：在自回归推理中，prefill 阶段为 prompt 建立每一层的 Key/Value 缓存，decode 阶段每生成一个新 token 都读取历史 KV，并追加新的 K/V。KV Cache 的基本显存公式为：

[
M_{\mathrm{KV}}
===============

2
\cdot
L
\cdot
B
\cdot
S
\cdot
n_{kv}
\cdot
d_h
\cdot
s_{\mathrm{dtype}}.
]

这个公式描述的是“真实有效 token”需要的理论显存。但在线 serving 系统中的难点并不只是理论显存，而是\textbf{如何在动态请求中高效分配、增长、复用和释放 KV Cache}。

如果所有请求长度相同、同时开始、同时结束，那么 KV Cache 管理会非常简单。系统只需要分配一个规则的矩形张量：

[
[L,B,n_{kv},S,d_h].
]

然而真实 LLM serving 中，请求具有高度动态性：

\begin{itemize}
\item 每个请求的 prompt length 不同；
\item 每个请求的 output length 不同；
\item 请求随时到达；
\item 请求随时结束；
\item 有些请求会被取消；
\item 有些请求会进入多轮对话；
\item 有些请求共享 prefix，有些请求完全不同；
\item 有些请求只生成几个 token，有些请求生成几千 token。
\end{itemize}

因此，KV Cache 更像一个在线内存分配系统，而不是一个静态张量。PagedAttention 的核心价值，正是把 KV Cache 管理从“连续大块分配”改造成“分页式 block 管理”，从而显著减少碎片和 over-reservation。

\subsection{不同请求长度导致的浪费}
\label{subsec:waste-from-variable-request-lengths}

考虑一个简单场景：系统同时服务 $B$ 个请求，每个请求最多允许上下文长度为 $S_{\max}$。若使用最朴素的连续分配方式，系统可能为每个请求都预留 $S_{\max}$ 个 token 的 KV Cache。

第 $i$ 个请求实际使用长度为：

[
S_i=S_{\mathrm{prompt},i}+S_{\mathrm{generated},i}.
]

如果为每个请求预留 $S_{\max}$，则第 $i$ 个请求的 cache 利用率为：

[
u_i=\frac{S_i}{S_{\max}}.
]

整个 batch 的平均利用率为：

[
U
=

\frac{\sum_{i=1}^{B}S_i}
{B S_{\max}}.
]

当请求长度差异很大时，$U$ 可能非常低。例如，一个 batch 中有如下请求：

[
[128,\ 256,\ 512,\ 2048,\ 4096]
]

若都按 $4096$ 预留，则总预留 token slots 为：

[
5\times 4096=20480.
]

真实使用 token slots 为：

[
128+256+512+2048+4096=7040.
]

利用率为：

[
U=\frac{7040}{20480}\approx 34.4%.
]

也就是说，超过一半的 KV Cache 显存只是被预留，却没有存放有效 token。

这类浪费在在线服务中非常常见。很多用户请求很短，例如“解释一个概念”“改写一句话”“写一个 SQL”；也有一些请求很长，例如“总结一篇长文档”“分析代码仓库”“多轮对话历史很长”。如果系统为了兼容长请求而给所有请求预留最大长度，短请求会造成大量内部浪费。

\begin{figure}[H]
\centering
\begin{tikzpicture}[
font=\small,
cell/.style={draw, minimum width=0.42cm, minimum height=0.38cm},
label/.style={font=\scriptsize},
used/.style={fill=blue!12},
unused/.style={fill=gray!18}
]

```
\node[font=\bfseries] at (0,3.05) {连续 KV Cache 预留导致的空间浪费};

\foreach \r/\y/\usedcells in {1/1.65/2,2/1.05/4,3/0.45/6,4/-0.15/10,5/-0.75/16} {
  \node[label] at (-4.8,\y) {Req \r};
  \foreach \i in {0,...,15} {
    \ifnum\i<\usedcells
      \node[cell, used] at ({-3.9+0.45*\i},\y) {};
    \else
      \node[cell, unused] at ({-3.9+0.45*\i},\y) {};
    \fi
  }
}

\node[draw, rounded corners=2pt, fill=blue!12, minimum width=0.55cm, minimum height=0.35cm] at (3.8,1.45) {};
\node[anchor=west] at (4.2,1.45) {used KV};
\node[draw, rounded corners=2pt, fill=gray!18, minimum width=0.55cm, minimum height=0.35cm] at (3.8,0.85) {};
\node[anchor=west] at (4.2,0.85) {reserved unused};

\node[align=left, text width=10.5cm] at (0,-2.0) {
  所有请求都按最大长度预留时，短请求会占用大量未使用 slots。
  对长短请求混部的在线 serving，这种 over-reservation 会显著降低可并发请求数。
};
```

\end{tikzpicture}
\caption{传统连续 KV Cache 分配中的长度浪费}
\label{fig:traditional-contiguous-fragmentation}
\end{figure}

\subsection{over-reservation}
\label{subsec:over-reservation}

Over-reservation 指系统为了避免未来扩容困难，提前为请求预留超过当前需要的 KV Cache 空间。它通常来自两个考虑：

\begin{itemize}
\item 请求未来还会继续 decode，cache 会增长；
\item 连续分配难以原地扩展，因此一开始预留最大长度更简单。
\end{itemize}

设请求当前长度为 $S_{\mathrm{cur}}$，最大允许长度为 $S_{\max}$。连续分配通常分配：

[
A_{\mathrm{contiguous}}=S_{\max}.
]

真实当前使用：

[
U_{\mathrm{cur}}=S_{\mathrm{cur}}.
]

当前浪费：

[
W_{\mathrm{over}}
=================

S_{\max}-S_{\mathrm{cur}}.
]

如果请求最终只生成到 $S_{\mathrm{final}}$，则最终浪费为：

[
W_{\mathrm{final}}
==================

S_{\max}-S_{\mathrm{final}}.
]

这类浪费尤其影响 admission control。假设 GPU 上剩余 KV Cache token capacity 为 $C_{\mathrm{free}}$。如果系统按最大长度预留，则新请求能否进入取决于：

[
S_{\max}\le C_{\mathrm{free}}.
]

但如果按实际增长分配，进入条件可能只需要：

[
S_{\mathrm{prompt}}+\Delta \le C_{\mathrm{free}},
]

其中 $\Delta$ 是少量初始 decode 余量。连续预留会使系统过早拒绝请求，降低吞吐。

Over-reservation 的本质是用显存换实现简单性。对于小模型或低并发服务，这可能可以接受；对于长上下文、高并发、GPU 显存昂贵的 LLM serving，它会成为主要成本来源。

\begin{table}[H]
\centering
\small
\caption{连续预留与按需分配的差异}
\label{tab:over-reservation-vs-demand-allocation}
\begin{tabular}{p{0.24\textwidth}p{0.34\textwidth}p{0.32\textwidth}}
\toprule
\textbf{维度} & \textbf{连续最大长度预留} & \textbf{按需增长分配} \
\midrule
实现复杂度
& 简单
& 需要动态内存管理 \
\midrule
访存模式
& 连续，kernel 读取简单
& 需要地址映射或 block table \
\midrule
显存利用率
& 长短请求混部时较低
& 通常更高 \
\midrule
请求扩容
& 无需扩容，已预留
& 生成过程中申请新 block \
\midrule
碎片风险
& 内部浪费严重，外部碎片也可能出现
& 主要是最后一个 block 的内部碎片 \
\midrule
适用场景
& 固定长度、离线批处理、低并发
& 在线 serving、长短请求混部、高并发 \
\bottomrule
\end{tabular}
\end{table}

\subsection{internal fragmentation}
\label{subsec:internal-fragmentation}

Internal fragmentation 指已经分配给某个请求的内存块内部，有一部分没有被使用。连续预留中的内部碎片可能非常大，因为每个请求都预留到最大长度。

PagedAttention 使用固定大小 block 后，仍然存在内部碎片，但碎片被限制在最后一个 block 内。

设 block size 为 $B_{\mathrm{blk}}$，请求实际长度为 $S_i$。需要分配的 block 数为：

[
N_{\mathrm{blk},i}
==================

\left\lceil
\frac{S_i}{B_{\mathrm{blk}}}
\right\rceil.
]

分配的 token slots 为：

[
A_i
===

N_{\mathrm{blk},i}B_{\mathrm{blk}}.
]

内部碎片为：

[
W_i
===

A_i-S_i.
]

因为只可能最后一个 block 没用满，所以：

[
0\le W_i < B_{\mathrm{blk}}.
]

这意味着每个请求最多浪费 $B_{\mathrm{blk}}-1$ 个 token slots。若 block size 为 $16$，每个请求最多浪费 $15$ 个 token slots；若 block size 为 $32$，最多浪费 $31$ 个 token slots。

与按 $S_{\max}$ 连续预留相比，这个浪费上界小得多。

[
W_{\mathrm{paged}}
<
B_{\mathrm{blk}},
]
[
W_{\mathrm{contiguous}}
=======================

S_{\max}-S_i.
]

当 $S_{\max}$ 为几千或几万，而 $B_{\mathrm{blk}}$ 只有几十时，分页式分配的空间效率优势非常明显。

\begin{figure}[H]
\centering
\begin{tikzpicture}[
font=\small,
cell/.style={draw, minimum width=0.38cm, minimum height=0.34cm},
blockbox/.style={draw, rounded corners=2pt, dashed, thick},
used/.style={fill=green!12},
unused/.style={fill=gray!18}
]

```
\node[font=\bfseries] at (0,3.0) {Block-based 分配中的内部碎片只出现在最后一个 block};

\node[anchor=east] at (-4.7,1.35) {logical tokens};
\foreach \i in {0,...,18} {
  \node[cell, used] at ({-4.1+0.4*\i},1.35) {};
}

\foreach \i in {19,...,23} {
  \node[cell, unused] at ({-4.1+0.4*\i},1.35) {};
}

\draw[blockbox] (-4.32,1.08) rectangle (-1.02,1.62);
\draw[blockbox] (-0.92,1.08) rectangle (2.38,1.62);
\draw[blockbox] (2.48,1.08) rectangle (5.78,1.62);

\node[font=\scriptsize] at (-2.67,0.75) {block 0 full};
\node[font=\scriptsize] at (0.73,0.75) {block 1 full};
\node[font=\scriptsize] at (4.13,0.75) {last block partially used};

\node[draw, rounded corners=2pt, fill=green!12, minimum width=0.55cm, minimum height=0.35cm] at (-1.0,-0.15) {};
\node[anchor=west] at (-0.6,-0.15) {used};
\node[draw, rounded corners=2pt, fill=gray!18, minimum width=0.55cm, minimum height=0.35cm] at (1.2,-0.15) {};
\node[anchor=west] at (1.6,-0.15) {internal fragmentation};

\node[align=left, text width=10.4cm] at (0,-1.45) {
  请求长度不是 block size 的整数倍时，最后一个 block 会有少量未使用 slots。
  这种浪费上界由 block size 决定，而不是由最大上下文长度决定。
};
```

\end{tikzpicture}
\caption{Block-based allocation 中的内部碎片}
\label{fig:block-internal-fragmentation}
\end{figure}

\subsection{为什么传统 batch 管理低效}
\label{subsec:why-traditional-batch-management-is-inefficient}

传统 batch 管理通常假设一个 batch 内的请求同步执行。对于普通深度学习推理，这很自然：一批输入进入模型，一次 forward 后全部完成。但 LLM decode 是迭代式的，一个请求可能生成 5 个 token，另一个请求可能生成 500 个 token。如果仍按静态 batch 管理，会出现严重低效。

静态 batch 的典型问题包括：

\begin{itemize}
\item \textbf{padding waste}：为了对齐长度，短请求被 padding 到长请求长度；
\item \textbf{head-of-line blocking}：短请求已经结束，但必须等长请求完成才能释放 batch；
\item \textbf{GPU 利用率波动}：batch 内请求逐渐结束，实际活跃请求数下降；
\item \textbf{KV Cache 释放不及时}：结束请求的 cache 不能立刻被新请求复用；
\item \textbf{新请求无法及时加入}：需要等整个 batch 结束；
\item \textbf{长请求拖慢短请求}：短请求 tail latency 被长请求影响。
\end{itemize}

Continuous batching 的思想是 iteration-level scheduling：每个 decode iteration 都可以让新请求加入，让完成请求退出。PagedAttention 提供了内存管理基础，使这种动态调度可行。因为请求不再绑定到固定连续 slot，而是持有一组可动态增长和释放的 blocks。

\begin{figure}[H]
\centering
\begin{tikzpicture}[
font=\small,
cell/.style={draw, minimum width=0.58cm, minimum height=0.42cm, align=center},
active/.style={fill=blue!12},
done/.style={fill=gray!18},
wait/.style={fill=orange!14}
]

```
\node[font=\bfseries] at (0,3.05) {Static Batching 中短请求等待长请求};

\foreach \t in {0,...,9} {
  \node[font=\scriptsize] at ({-3.7+0.65*\t},2.25) {t\t};
}

\foreach \r/\y/\len in {1/1.55/3,2/0.95/5,3/0.35/9} {
  \node[font=\scriptsize] at (-4.55,\y) {Req \r};
  \foreach \t in {0,...,9} {
    \ifnum\t<\len
      \node[cell, active] at ({-3.7+0.65*\t},\y) {run};
    \else
      \node[cell, done] at ({-3.7+0.65*\t},\y) {};
    \fi
  }
}

\node[cell, wait, minimum width=1.6cm] at (3.85,0.95) {new req waits};

\node[align=left, text width=10.4cm] at (0,-1.0) {
  静态 batch 中，新请求通常要等当前 batch 结束才能加入。
  短请求虽然已经完成，但资源释放和新请求调度都被长请求阻塞。
};
```

\end{tikzpicture}
\caption{传统静态 batch 管理中的 head-of-line blocking}
\label{fig:static-batching-head-of-line-blocking}
\end{figure}

\section{PagedAttention 核心思想}
\label{sec:paged-attention-core-idea}

PagedAttention 的名字来自操作系统中的分页内存管理。它不是改变 attention 的数学公式，而是改变 KV Cache 的\textbf{物理存储方式}和 attention kernel 的\textbf{读取方式}。

标准 attention 仍然是：

[
\operatorname{Attention}(Q,K,V)
===============================

\softmax
\left(
\frac{QK^{\mathsf{T}}}{\sqrt{d_h}}
\right)V.
]

PagedAttention 改变的是：逻辑上连续的 $K,V$ 序列，不再要求物理上连续存储。每个请求的逻辑 token 序列被切成多个 logical blocks；系统从全局 physical block pool 中分配 physical blocks；block table 记录 logical block 到 physical block 的映射。Attention kernel 在读取历史 K/V 时，通过 block table 做地址翻译。

\subsection{操作系统分页类比}
\label{subsec:os-paging-analogy}

操作系统虚拟内存中，进程看到的是连续虚拟地址空间。实际物理内存可能是不连续的 page frames。页表 page table 负责把虚拟页号映射到物理页框。

PagedAttention 中，请求看到的是连续 token positions：

[
0,1,2,\ldots,S-1.
]

但这些 token 的 KV Cache 可以分散存放在不同 physical blocks 中。Block table 类似页表：

[
\text{logical block id}
\rightarrow
\text{physical block id}.
]

类比如下：

\begin{table}[H]
\centering
\small
\caption{操作系统分页与 PagedAttention 的类比}
\label{tab:os-paging-pagedattention-analogy}
\begin{tabular}{p{0.28\textwidth}p{0.30\textwidth}p{0.32\textwidth}}
\toprule
\textbf{概念} & \textbf{操作系统分页} & \textbf{PagedAttention} \
\midrule
地址空间
& 进程虚拟地址空间
& 请求的逻辑 token 序列 \
\midrule
页
& virtual page
& logical KV block \
\midrule
物理页框
& physical page frame
& physical KV block \
\midrule
页表
& page table
& block table \
\midrule
地址翻译
& virtual address $\rightarrow$ physical address
& token position $\rightarrow$ physical KV address \
\midrule
共享页
& 多进程共享只读页
& 多请求共享 prefix blocks \
\midrule
Copy-on-write
& 写共享页前复制
& 请求分叉或 beam search 时复制 block \
\bottomrule
\end{tabular}
\end{table}

这个类比非常重要，因为它解释了 PagedAttention 的设计哲学：不要要求每个请求在物理显存中连续，而是让 kernel 支持一次轻量级地址翻译。用少量寻址开销换取更高显存利用率和更灵活的调度。

\begin{figure}[H]
\centering
\begin{tikzpicture}[
font=\small,
box/.style={draw, rounded corners=3pt, minimum width=1.9cm, minimum height=0.65cm, align=center},
page/.style={draw, rounded corners=2pt, minimum width=0.68cm, minimum height=0.46cm, align=center},
arrow/.style={-{Latex[length=2.1mm]}, thick}
]

```
\node[font=\bfseries] at (0,3.25) {操作系统分页与 PagedAttention 的结构类比};

% OS side
\node[font=\bfseries] at (-3.7,2.45) {OS Paging};
\foreach \i/\x in {0/-5.0,1/-4.25,2/-3.5,3/-2.75} {
  \node[page, fill=blue!10] (vp\i) at (\x,1.75) {VP\i};
}
\node[box, fill=orange!14] (pt) at (-3.85,0.75) {Page Table};
\foreach \i/\x in {0/-5.0,1/-4.25,2/-3.5,3/-2.75} {
  \node[page, fill=green!10] (pf\i) at (\x,-0.35) {PF\i};
}
\foreach \i in {0,1,2,3} {
  \draw[arrow] (vp\i) -- (pt);
  \draw[arrow] (pt) -- (pf\i);
}

% PA side
\node[font=\bfseries] at (3.7,2.45) {PagedAttention};
\foreach \i/\x in {0/2.35,1/3.1,2/3.85,3/4.6} {
  \node[page, fill=blue!10] (lb\i) at (\x,1.75) {LB\i};
}
\node[box, fill=orange!14] (bt) at (3.5,0.75) {Block Table};
\foreach \i/\x/\pid in {0/2.35/7,1/3.1/2,2/3.85/9,3/4.6/5} {
  \node[page, fill=green!10] (pb\i) at (\x,-0.35) {PB\pid};
}
\foreach \i in {0,1,2,3} {
  \draw[arrow] (lb\i) -- (bt);
  \draw[arrow] (bt) -- (pb\i);
}

\draw[dashed, thick] (0,-0.9) -- (0,2.7);

\node[align=left, text width=10.5cm] at (0,-1.95) {
  请求逻辑上看到连续 token positions，但物理 KV blocks 可以分散存放。
  block table 扮演页表角色，把逻辑位置映射到物理显存地址。
};
```

\end{tikzpicture}
\caption{操作系统分页与 PagedAttention 的类比}
\label{fig:os-paging-vs-pagedattention}
\end{figure}

\subsection{logical block 与 physical block}
\label{subsec:logical-block-and-physical-block}

设 block size 为：

[
B_{\mathrm{blk}}.
]

请求的逻辑 token 位置 $s$ 可以分解为：

[
b_{\mathrm{logical}}
====================

\left\lfloor
\frac{s}{B_{\mathrm{blk}}}
\right\rfloor,
]
[
o
=

s\bmod B_{\mathrm{blk}}.
]

其中 $b_{\mathrm{logical}}$ 是 logical block id，$o$ 是 block 内 offset。

Block table 给出：

[
b_{\mathrm{physical}}
=====================

\operatorname{BlockTable}[r,b_{\mathrm{logical}}],
]

其中 $r$ 是 request id 或 sequence id。

最终物理地址可以抽象为：

[
\operatorname{addr}
===================

\operatorname{Base}
+
b_{\mathrm{physical}}
\cdot
\operatorname{BlockStride}
+
o
\cdot
\operatorname{TokenStride}
+
h
\cdot
\operatorname{HeadStride}
+
d
\cdot
\operatorname{DimStride}.
]

这里 $h$ 是 KV head index，$d$ 是 head dimension index。真实 kernel 会根据具体 layout 优化 stride 和向量化加载。

Logical block 和 physical block 的分离带来三个重要能力：

\begin{itemize}
\item \textbf{非连续物理存储}：请求的逻辑 blocks 可以映射到任意 physical blocks；
\item \textbf{按需增长}：请求需要更多 cache 时，只需追加新的 physical block；
\item \textbf{共享与复制}：多个请求可以引用相同 physical block，写入时再 copy-on-write。
\end{itemize}

\subsection{block table}
\label{subsec:block-table}

Block table 是 PagedAttention 的核心元数据。对于每个请求，它记录每个 logical block 对应的 physical block：

[
\operatorname{BlockTable}
\in
\mathbb{N}^{N_{\mathrm{req}}\times N_{\mathrm{max_blocks}}}.
]

例如，请求 $r$ 的 block table 为：

[
[7,\ 2,\ 9,\ 5].
]

这表示：

[
\text{logical block }0\rightarrow\text{physical block }7,
]
[
\text{logical block }1\rightarrow\text{physical block }2,
]
[
\text{logical block }2\rightarrow\text{physical block }9,
]
[
\text{logical block }3\rightarrow\text{physical block }5.
]

请求逻辑序列仍然是连续的，但 physical blocks 可以任意分散。

\begin{figure}[H]
\centering
\begin{tikzpicture}[
font=\small,
lb/.style={draw, rounded corners=2pt, minimum width=0.85cm, minimum height=0.48cm, align=center, fill=blue!10},
pb/.style={draw, rounded corners=2pt, minimum width=0.85cm, minimum height=0.48cm, align=center, fill=green!10},
table/.style={draw, rounded corners=3pt, minimum width=2.5cm, minimum height=1.3cm, align=center, fill=orange!14},
arrow/.style={-{Latex[length=2.1mm]}, thick}
]

```
\node[font=\bfseries] at (0,3.25) {PagedAttention block table 映射};

\foreach \i/\x in {0/-5.0,1/-4.0,2/-3.0,3/-2.0} {
  \node[lb] (lb\i) at (\x,1.65) {LB \i};
}

\node[table] (bt) at (-3.5,0.45) {
  Block Table\\
  \begin{tabular}{c|c}
  LB & PB\\
  0 & 7\\
  1 & 2\\
  2 & 9\\
  3 & 5
  \end{tabular}
};

\foreach \i/\x/\pid in {0/1.2/2,1/2.2/5,2/3.2/7,3/4.2/9} {
  \node[pb] (pb\pid) at (\x,1.65) {PB \pid};
}

\foreach \i in {0,1,2,3} {
  \draw[arrow] (lb\i) -- (bt);
}

\draw[arrow] (bt) -- (pb7);
\draw[arrow] (bt) -- (pb2);
\draw[arrow] (bt) -- (pb9);
\draw[arrow] (bt) -- (pb5);

\node[align=left, text width=10.4cm] at (0,-1.35) {
  block table 把请求的逻辑 block 映射到全局 physical block pool 中的任意 block。
  Attention kernel 根据该表读取逻辑上连续、物理上不连续的 KV Cache。
};
```

\end{tikzpicture}
\caption{PagedAttention 中 logical block 到 physical block 的映射}
\label{fig:pagedattention-block-table}
\end{figure}

Block table 的大小通常远小于 KV Cache 本身。若每个 block 存 $B_{\mathrm{blk}}$ 个 token，一个请求长度为 $S$，则 block table 只需要：

[
\left\lceil\frac{S}{B_{\mathrm{blk}}}\right\rceil
]

个整数。相比每个 token 的 K/V 向量，metadata 开销很小。

\subsection{attention kernel 如何读取分页 KV}
\label{subsec:how-attention-kernel-reads-paged-kv}

PagedAttention kernel 的核心是在遍历历史 token 时进行地址翻译。对于当前 query token，attention 需要遍历历史位置：

[
s=0,1,\ldots,S-1.
]

对于每个 $s$，kernel 计算：

[
b=\left\lfloor\frac{s}{B_{\mathrm{blk}}}\right\rfloor,
\quad
o=s\bmod B_{\mathrm{blk}}.
]

然后从 block table 读取 physical block id：

[
p=\operatorname{BlockTable}[r,b].
]

再从物理 KV Cache 中读取：

[
\mat{K}*{phys}[p,o,h,:],
\quad
\mat{V}*{phys}[p,o,h,:].
]

因此，PagedAttention 的读取路径可以写成：

[
s
\rightarrow
(b,o)
\rightarrow
p
\rightarrow
\operatorname{addr}(p,o,h,d)
\rightarrow
K/V.
]

下面是一个高度简化的伪代码，用于说明逻辑，不代表高性能 kernel。

\begin{lstlisting}[language=Python,caption={PagedAttention 中根据 block table 读取 KV 的伪代码},label={lst:pagedattention-read-kv-pseudocode}]
def read_kv_for_position(
request_id,
position,
kv_head,
block_size,
block_table,
k_blocks,
v_blocks,
):
logical_block = position // block_size
offset = position % block_size

```
physical_block = block_table[request_id][logical_block]

# k_blocks/v_blocks:
# [num_physical_blocks, num_kv_heads, block_size, head_dim]
k = k_blocks[physical_block, kv_head, offset, :]
v = v_blocks[physical_block, kv_head, offset, :]

return k, v
```

\end{lstlisting}

真正的 PagedAttention kernel 需要处理：

\begin{itemize}
\item 多个 requests；
\item 多个 heads；
\item 每个请求不同 context length；
\item 向量化加载 K/V；
\item block 内连续访存；
\item softmax 的数值稳定；
\item warp/block 级并行；
\item shared memory 与 register 使用；
\item GQA/MQA head mapping；
\item causal mask；
\item 最后一个 block 未满的边界；
\item tensor parallel 下的 head shard。
\end{itemize}

PagedAttention 的挑战在于：它必须在引入 block table 间接寻址的同时，保持 attention kernel 的访存和并行效率。若地址翻译开销过大，显存节省可能被 kernel 开销抵消。因此 block size、layout、kernel tiling 和调度策略必须共同设计。

\begin{figure}[H]
\centering
\begin{tikzpicture}[
font=\small,
box/.style={draw, rounded corners=3pt, minimum width=1.75cm, minimum height=0.65cm, align=center},
arrow/.style={-{Latex[length=2.1mm]}, thick}
]

```
\node[font=\bfseries] at (0,3.25) {PagedAttention kernel 的 KV 读取路径};

\node[box, fill=blue!6] (pos) at (-5.2,1.1) {logical\\position $s$};
\node[box, fill=orange!14] (split) at (-3.0,1.1) {$s\rightarrow(b,o)$};
\node[box, fill=green!10] (table) at (-0.75,1.1) {Block\\Table};
\node[box, fill=yellow!16] (phys) at (1.55,1.1) {physical\\block $p$};
\node[box, fill=purple!10] (addr) at (3.8,1.1) {KV memory\\address};
\node[box, fill=red!8] (kv) at (5.85,1.1) {$K,V$};

\draw[arrow] (pos) -- (split);
\draw[arrow] (split) -- (table);
\draw[arrow] (table) -- (phys);
\draw[arrow] (phys) -- (addr);
\draw[arrow] (addr) -- (kv);

\node[align=left, text width=10.6cm] at (0,-1.0) {
  PagedAttention 在 attention kernel 内部完成逻辑位置到物理 block 的地址翻译。
  逻辑序列连续，物理存储可以离散，从而避免大块连续 KV Cache 分配。
};
```

\end{tikzpicture}
\caption{PagedAttention kernel 通过 block table 读取分页 KV}
\label{fig:pagedattention-kernel-read-path}
\end{figure}

\subsection{copy-on-write 与共享 block}
\label{subsec:copy-on-write-shared-blocks}

PagedAttention 也天然支持共享 prefix。多个请求如果共享相同前缀，可以让它们的 block table 指向相同 physical blocks。每个 physical block 维护引用计数：

[
\operatorname{refcount}(p).
]

当多个请求共享 block 时：

[
\operatorname{refcount}(p)>1.
]

如果某个请求需要写入一个共享 block，而该写入会改变 block 内容，就不能直接原地修改，否则会影响其他请求。此时需要 copy-on-write：

[
p_{\mathrm{new}}=\operatorname{AllocateBlock}(),
]
[
\operatorname{Copy}(p_{\mathrm{old}},p_{\mathrm{new}}),
]
[
\operatorname{BlockTable}[r,b]\leftarrow p_{\mathrm{new}}.
]

在 LLM decode 中，通常只会向最后一个 block 追加 token。共享 prefix blocks 大多是只读的；分叉、beam search、parallel sampling 或 speculative decoding 可能会产生共享后再分叉的场景，此时 copy-on-write 很有用。

\begin{lstlisting}[language=Python,caption={Block 引用计数与 copy-on-write 的简化伪代码},label={lst:block-copy-on-write-pseudocode}]
def ensure_writable_block(request, logical_block_id, block_manager):
physical_block = request.block_table[logical_block_id]

```
if block_manager.ref_count[physical_block] == 1:
    return physical_block

# Copy-on-write.
new_block = block_manager.allocate_block()
copy_block(src=physical_block, dst=new_block)

block_manager.ref_count[physical_block] -= 1
block_manager.ref_count[new_block] = 1

request.block_table[logical_block_id] = new_block
return new_block
```

\end{lstlisting}

Copy-on-write 的意义不只是节省显存，也让复杂 decoding 策略更容易实现。例如 beam search 中多个 beams 共享同一个 prompt prefix，直到它们生成不同 suffix。若每个 beam 都复制完整 KV Cache，显存会迅速膨胀；若共享 prefix blocks，只复制分叉后的 blocks，就能显著降低开销。

\section{vLLM 架构}
\label{sec:vllm-architecture}

vLLM 是围绕高吞吐 LLM serving 设计的推理系统。它的核心思想之一是使用 PagedAttention 管理 KV Cache，并配合 continuous batching、scheduler、worker、block manager 等组件，让在线请求能动态加入、退出和共享 cache。

本节不追逐某个版本的具体代码结构，而从系统架构角度解释一个 vLLM 风格 serving engine 通常如何组织。

\subsection{总体架构}
\label{subsec:vllm-overall-architecture}

一个 vLLM 风格的 serving 系统可以抽象为以下组件：

\begin{itemize}
\item \textbf{API Server / Frontend}：接收 HTTP/gRPC/OpenAI-compatible 请求，处理鉴权、限流、请求解析和流式输出；
\item \textbf{Tokenizer}：将文本转为 token ids，将生成 token 转回文本；
\item \textbf{Scheduler}：决定每个 iteration 哪些请求运行 prefill，哪些请求运行 decode，哪些请求等待；
\item \textbf{Block Manager}：管理 KV Cache physical blocks 的分配、释放、引用计数和 block table；
\item \textbf{Worker / Executor}：在 GPU 上执行模型 forward、PagedAttention kernel、sampling 等；
\item \textbf{Model Runner}：封装模型权重、attention backend、KV Cache layout 和 CUDA graph 等执行细节；
\item \textbf{Output Processor}：处理采样、stop condition、流式 token 返回和请求结束逻辑。
\end{itemize}

它们的关系如下图所示。

\begin{figure}[H]
\centering
\begin{tikzpicture}[
font=\small,
comp/.style={draw, rounded corners=4pt, minimum width=2.05cm, minimum height=0.75cm, align=center},
gpu/.style={draw, rounded corners=4pt, minimum width=2.15cm, minimum height=0.85cm, align=center, fill=blue!6},
arrow/.style={-{Latex[length=2.2mm]}, thick}
]

```
\node[font=\bfseries] at (0,3.5) {vLLM 风格 Serving 架构};

\node[comp, fill=orange!12] (api) at (-5.2,1.95) {API Server\\Frontend};
\node[comp, fill=green!10] (tok) at (-2.75,1.95) {Tokenizer};
\node[comp, fill=yellow!16] (sched) at (-0.25,1.95) {Scheduler};
\node[comp, fill=purple!10] (bm) at (2.35,1.95) {Block\\Manager};
\node[gpu] (worker) at (4.95,1.95) {GPU Worker\\Model Runner};

\node[comp, fill=red!8] (out) at (-0.25,0.35) {Output\\Processor};
\node[comp, fill=blue!8] (kv) at (2.35,0.35) {KV Blocks\\Block Tables};
\node[gpu] (kernels) at (4.95,0.35) {PagedAttention\\Sampling Kernels};

\draw[arrow] (api) -- (tok);
\draw[arrow] (tok) -- (sched);
\draw[arrow] (sched) -- (bm);
\draw[arrow] (bm) -- (worker);
\draw[arrow] (worker) -- (kernels);
\draw[arrow] (bm) -- (kv);
\draw[arrow] (kv) -- (kernels);
\draw[arrow] (kernels) -- (out);
\draw[arrow] (out.west) .. controls (-2.2,0.15) and (-4.6,0.6) .. (api.south);

\node[align=left, text width=10.6cm] at (0,-1.25) {
  Scheduler 决定每个 iteration 的请求集合；Block Manager 为请求分配 KV blocks 并维护 block table；
  Worker 在 GPU 上执行模型和 PagedAttention kernel，Output Processor 处理采样结果与流式返回。
};
```

\end{tikzpicture}
\caption{vLLM 风格 serving engine 的核心组件}
\label{fig:vllm-serving-architecture}
\end{figure}

\subsection{scheduler}
\label{subsec:vllm-scheduler}

Scheduler 是 LLM serving 的大脑。它每个 iteration 都要决定：

\begin{itemize}
\item 哪些新请求可以进入系统；
\item 哪些请求执行 prefill；
\item 哪些请求执行 decode；
\item 哪些请求需要等待；
\item 哪些请求因为资源不足被延迟；
\item 哪些请求已经完成，需要释放资源；
\item 本 iteration 的 token budget 如何分配。
\end{itemize}

与传统 batch scheduler 不同，LLM scheduler 的基本调度单位不是“整个请求”，而是“本 iteration 的 token”。对于 decode 请求，每个 running request 通常只生成一个 token；对于 prefill 请求，一个 prompt 可能包含很多 token，系统可能一次处理完整 prompt，也可能使用 chunked prefill 将长 prompt 分块处理。

Scheduler 常维护几类队列：

\begin{itemize}
\item \textbf{waiting queue}：新到达但尚未分配资源的请求；
\item \textbf{running set}：已经拥有 KV blocks、正在 decode 或 prefill 的请求；
\item \textbf{finished set}：已完成，需要释放 blocks 并返回结果；
\item \textbf{paused/swapped set}：资源压力下暂时暂停或迁出的请求。
\end{itemize}

一个简化调度目标可以写成：

[
\max
\quad
\operatorname{tokens_scheduled}
]

同时满足：

[
\operatorname{KVBlocksNeeded}
\le
\operatorname{FreeBlocks},
]

[
\operatorname{BatchedTokens}
\le
T_{\max},
]

[
\operatorname{NumSequences}
\le
N_{\max}.
]

其中 $T_{\max}$ 是每 iteration 最大 token budget，$N_{\max}$ 是最大并发序列数。调度器需要在吞吐和延迟之间折中：多塞请求可以提升吞吐，但可能增加单请求延迟；优先短请求可以改善平均延迟，但可能让长请求饥饿；优先 decode 可以降低 TPOT，但可能让新请求 TTFT 变差。

\begin{lstlisting}[language=Python,caption={Iteration-level scheduler 的简化伪代码},label={lst:vllm-scheduler-pseudocode}]
def schedule_iteration(waiting, running, block_manager, token_budget, max_seqs):
scheduled_prefills = []
scheduled_decodes = []

```
remaining_tokens = token_budget
remaining_slots = max_seqs - len(running)

# 1. Schedule decode for already running requests.
for req in list(running):
    if req.is_finished():
        continue

    if not block_manager.can_append_slot(req):
        continue

    scheduled_decodes.append(req)
    remaining_tokens -= 1

    if remaining_tokens <= 0:
        break

# 2. Admit new requests if there is budget and KV memory.
while waiting and remaining_tokens > 0 and remaining_slots > 0:
    req = waiting[0]
    prefill_tokens = req.prompt_len

    if prefill_tokens > remaining_tokens:
        break

    if not block_manager.can_allocate_for_prefill(req):
        break

    waiting.pop(0)
    block_manager.allocate_for_prefill(req)
    running.add(req)

    scheduled_prefills.append(req)
    remaining_tokens -= prefill_tokens
    remaining_slots -= 1

return scheduled_prefills, scheduled_decodes
```

\end{lstlisting}

真实 scheduler 会更复杂。它会考虑 chunked prefill、prefix cache 命中、优先级、超时、beam search、多模态输入、LoRA adapter、speculative decoding、抢占和资源回收。但核心仍然是：在每个 iteration 中，根据 KV Cache 资源和 token budget 选择请求集合。

\subsection{worker}
\label{subsec:vllm-worker}

Worker 负责在 GPU 上执行实际计算。它通常接收 scheduler 产生的执行计划，包括：

\begin{itemize}
\item 本次 prefill 的 requests；
\item 本次 decode 的 requests；
\item input token ids；
\item position ids；
\item slot mapping；
\item block tables；
\item sequence lengths；
\item sampling parameters；
\item LoRA 或其他 adapter 信息。
\end{itemize}

Worker 的执行流程可以抽象为：

\begin{enumerate}
\item 根据调度计划准备输入张量；
\item 调用模型 forward；
\item attention kernel 根据 block table 读取 KV Cache；
\item 将新 token 的 K/V 写入对应 physical blocks；
\item 计算 logits；
\item 执行 sampling；
\item 返回 sampled tokens 和请求状态。
\end{enumerate}

Worker 的关键性能目标是减少 CPU/GPU 同步、减少小 kernel、减少动态 shape 造成的开销，并尽量让 GPU 一直有足够工作。

在高性能实现中，worker 可能使用 CUDA graph 固化部分执行路径，以减少 launch overhead。但 LLM serving 的动态性很强，batch size、sequence length、请求组合都在变化，CUDA graph 捕获和复用需要做 shape bucketing 或其他工程折中。

\subsection{block manager}
\label{subsec:vllm-block-manager}

Block Manager 是 PagedAttention 系统中的内存管理器。它维护全局 physical block pool，并为每个请求维护 block table。

Block Manager 的职责包括：

\begin{itemize}
\item 初始化 GPU KV block pool；
\item 为 prefill 请求分配足够 blocks；
\item 为 decode 追加 token 时分配新 block；
\item 维护 request 到 block table 的映射；
\item 维护 physical block 引用计数；
\item 请求结束时释放 blocks；
\item prefix cache 命中时共享 blocks；
\item copy-on-write；
\item 在资源不足时触发 eviction、pause 或拒绝请求。
\end{itemize}

一个极简 Block Manager 可以写成：

\begin{lstlisting}[language=Python,caption={Block Manager 的简化实现骨架},label={lst:block-manager-skeleton}]
import math
from collections import deque

class BlockManager:
def **init**(self, num_blocks, block_size):
self.num_blocks = num_blocks
self.block_size = block_size

```
    self.free_blocks = deque(range(num_blocks))
    self.ref_count = [0 for _ in range(num_blocks)]

    # request_id -> list of physical block ids
    self.block_tables = {}

def num_required_blocks(self, num_tokens):
    return math.ceil(num_tokens / self.block_size)

def can_allocate(self, num_tokens):
    return len(self.free_blocks) >= self.num_required_blocks(num_tokens)

def allocate(self, request_id, num_tokens):
    n = self.num_required_blocks(num_tokens)
    if len(self.free_blocks) < n:
        raise RuntimeError("not enough KV blocks")

    blocks = []
    for _ in range(n):
        block = self.free_blocks.popleft()
        self.ref_count[block] = 1
        blocks.append(block)

    self.block_tables[request_id] = blocks
    return blocks

def append_slot(self, request_id, current_length):
    blocks = self.block_tables[request_id]

    # Allocate a new block when current length reaches a block boundary.
    if current_length % self.block_size == 0:
        if not self.free_blocks:
            raise RuntimeError("no free block for decode")
        block = self.free_blocks.popleft()
        self.ref_count[block] = 1
        blocks.append(block)

    return blocks

def release(self, request_id):
    blocks = self.block_tables.pop(request_id, [])

    for block in blocks:
        self.ref_count[block] -= 1
        if self.ref_count[block] == 0:
            self.free_blocks.append(block)
```

\end{lstlisting}

真实 Block Manager 还要处理共享 prefix、copy-on-write、CPU swap、并发安全、GPU/CPU metadata 同步和调度器交互。它是 serving 系统中最容易影响稳定性和显存利用率的组件之一。

\subsection{continuous batching}
\label{subsec:vllm-continuous-batching}

Continuous batching 指系统在每个 decode iteration 重新组织 batch，让新请求可以加入，让完成请求可以退出。它与 PagedAttention 是互补关系：

\begin{itemize}
\item Continuous batching 解决\textbf{请求调度动态性}；
\item PagedAttention 解决\textbf{动态请求的 KV Cache 内存管理}。
\end{itemize}

如果没有 PagedAttention，continuous batching 仍然会遇到连续 KV Cache 分配困难：请求加入和退出会导致 batch slot 移动、cache 拷贝、碎片和 over-reservation。PagedAttention 让请求的 cache 独立于 batch slot，只需要更新 block table 和调度 metadata。

一个 iteration-level timeline 如下：

\begin{figure}[H]
\centering
\begin{tikzpicture}[
font=\scriptsize,
cell/.style={draw, minimum width=0.7cm, minimum height=0.42cm, align=center},
decode/.style={fill=blue!12},
prefill/.style={fill=green!12},
done/.style={fill=gray!18},
newreq/.style={fill=orange!14}
]

```
\node[font=\bfseries\small] at (0,3.0) {vLLM 风格 Continuous Batching 时间线};

\foreach \t in {0,...,8} {
  \node at ({-3.6+0.78*\t},2.35) {iter \t};
}

\foreach \r/\y in {A/1.65,B/1.1,C/0.55,D/0.0} {
  \node at (-4.45,\y) {Req \r};
}

% A decode then finish
\foreach \t in {0,1,2,3} {
  \node[cell, decode] at ({-3.6+0.78*\t},1.65) {D};
}
\foreach \t in {4,...,8} {
  \node[cell, done] at ({-3.6+0.78*\t},1.65) {};
}

% B long decode
\foreach \t in {0,...,8} {
  \node[cell, decode] at ({-3.6+0.78*\t},1.1) {D};
}

% C joins at 2
\foreach \t in {0,1} {
  \node[cell, done] at ({-3.6+0.78*\t},0.55) {};
}
\node[cell, prefill] at ({-3.6+0.78*2},0.55) {P};
\foreach \t in {3,...,8} {
  \node[cell, decode] at ({-3.6+0.78*\t},0.55) {D};
}

% D joins at 5
\foreach \t in {0,...,4} {
  \node[cell, done] at ({-3.6+0.78*\t},0.0) {};
}
\node[cell, prefill] at ({-3.6+0.78*5},0.0) {P};
\foreach \t in {6,...,8} {
  \node[cell, decode] at ({-3.6+0.78*\t},0.0) {D};
}

\node[draw, rounded corners=2pt, fill=green!12, minimum width=0.5cm, minimum height=0.32cm] at (3.8,1.5) {};
\node[anchor=west] at (4.15,1.5) {Prefill};
\node[draw, rounded corners=2pt, fill=blue!12, minimum width=0.5cm, minimum height=0.32cm] at (3.8,1.0) {};
\node[anchor=west] at (4.15,1.0) {Decode};
\node[draw, rounded corners=2pt, fill=gray!18, minimum width=0.5cm, minimum height=0.32cm] at (3.8,0.5) {};
\node[anchor=west] at (4.15,0.5) {inactive};

\node[align=left, text width=10.5cm, font=\small] at (0,-1.35) {
  每个 iteration 都可以重新组成 batch。已完成请求释放 blocks，新请求分配 blocks 后加入。
  PagedAttention 让这种动态调度不需要移动已有请求的 KV Cache。
};
```

\end{tikzpicture}
\caption{Continuous batching 中请求动态加入和退出}
\label{fig:vllm-continuous-batching-timeline}
\end{figure}

图中 P 表示 prefill，D 表示 decode。Req C 和 Req D 在不同 iteration 加入；Req A 结束后释放资源；Req B 持续 decode。每个 iteration 的实际 batch 都不同。

\subsection{prefix caching 与 speculative decoding}
\label{subsec:vllm-prefix-caching-speculative-decoding}

Prefix caching 与 PagedAttention 结合得很自然，因为共享 prefix 可以通过共享 physical blocks 实现。多个请求的 block table 前缀指向相同 blocks，refcount 增加即可。只有当请求在共享 prefix 后继续生成自己的 suffix 时，才分配新的 blocks。

Speculative decoding 则引入另一类 cache 管理问题。它通常包含 draft model 和 target model：draft model 先快速生成多个候选 token，target model 验证这些 token。若候选 token 被接受，则可以一次提交多个 token；若被拒绝，需要回滚到某个位置。

从 KV Cache 管理角度看，speculative decoding 需要处理：

\begin{itemize}
\item 草稿 token 的临时 KV；
\item 被接受 token 的 KV 提交；
\item 被拒绝 token 的 KV 回滚；
\item draft 和 target 的 cache 是否分开；
\item 多 token 提交时 block table 的更新；
\item 回滚时释放未提交 blocks。
\end{itemize}

Paged block 管理对这些场景有帮助，因为未提交 token 可能只占用少量 blocks；回滚时释放这些 blocks 即可，而不需要移动连续大 cache。

\begin{figure}[H]
\centering
\begin{tikzpicture}[
font=\small,
block/.style={draw, rounded corners=2pt, minimum width=0.55cm, minimum height=0.38cm, align=center},
arrow/.style={-{Latex[length=2.0mm]}, thick}
]

```
\node[font=\bfseries] at (0,3.15) {Speculative Decoding 中的提交与回滚};

\foreach \i/\x in {0/-4.0,1/-3.4,2/-2.8,3/-2.2} {
  \node[block, fill=green!12] (c\i) at (\x,1.4) {$c$};
}
\node[anchor=east] at (-4.5,1.4) {committed};

\foreach \i/\x in {0/-1.2,1/-0.6,2/0.0,3/0.6} {
  \node[block, fill=orange!14] (d\i) at (\x,1.4) {$d$};
}
\node[anchor=south] at (-0.3,1.85) {draft tokens};

\foreach \i/\x in {0/1.8,1/2.4} {
  \node[block, fill=blue!12] at (\x,1.4) {$a$};
}
\foreach \i/\x in {0/3.0,1/3.6} {
  \node[block, fill=gray!18] at (\x,1.4) {$r$};
}

\draw[arrow] (0.9,0.9) -- node[below] {verify} (1.55,0.9);

\node[align=center, font=\scriptsize] at (2.1,0.65) {accepted};
\node[align=center, font=\scriptsize] at (3.35,0.65) {rollback};

\node[align=left, text width=10.5cm] at (0,-0.95) {
  Speculative decoding 会产生临时 token。验证通过的 token 提交到正式 KV Cache；
  未通过的 token 对应 blocks 需要释放或回滚。分页式 block 管理让这种操作更自然。
};
```

\end{tikzpicture}
\caption{Speculative decoding 中 KV blocks 的提交与回滚}
\label{fig:speculative-decoding-kv-commit-rollback}
\end{figure}

\section{PagedAttention 的性能收益}
\label{sec:pagedattention-performance-benefits}

PagedAttention 的收益不能简单理解为“attention 算法更少 FLOPs”。它的主要价值在于\textbf{提高 KV Cache 显存利用率}，从而允许更多请求同时运行，并使 continuous batching 更有效。吞吐提升往往来自更高并发、更低内存浪费和更好的调度，而不是单个 attention step 的数学计算减少。

\subsection{显存利用率提升}
\label{subsec:pagedattention-memory-utilization}

连续分配中，每个请求分配 $S_{\max}$ token slots：

[
A_{\mathrm{contiguous},i}=S_{\max}.
]

PagedAttention 中，每个请求按 block 分配：

[
A_{\mathrm{paged},i}
====================

B_{\mathrm{blk}}
\left\lceil
\frac{S_i}{B_{\mathrm{blk}}}
\right\rceil.
]

总利用率分别为：

[
U_{\mathrm{contiguous}}
=======================

\frac{\sum_i S_i}
{\sum_i S_{\max}},
]

[
U_{\mathrm{paged}}
==================

\frac{\sum_i S_i}
{\sum_i B_{\mathrm{blk}}
\left\lceil
\frac{S_i}{B_{\mathrm{blk}}}
\right\rceil}.
]

当 $B_{\mathrm{blk}}\ll S_{\max}$ 时，PagedAttention 的利用率通常显著更高。

下面用一个简单例子说明。假设 block size 为 $16$，请求长度为：

[
[128,\ 256,\ 512,\ 2048,\ 4096].
]

PagedAttention 分配 slots 为：

[
[128,\ 256,\ 512,\ 2048,\ 4096],
]

因为这些长度恰好是 16 的倍数，利用率接近 $100%$。若长度为：

[
[130,\ 257,\ 519,\ 2050,\ 4097],
]

则分配 slots 为：

[
[144,\ 272,\ 528,\ 2064,\ 4112].
]

总使用为：

[
7053.
]

总分配为：

[
7120.
]

利用率为：

[
\frac{7053}{7120}\approx 99.1%.
]

相比连续最大长度预留的约 $34%$，差异非常大。当然真实系统中还有 metadata、block table、alignment 和其他 buffer，但核心趋势成立。

\begin{table}[H]
\centering
\small
\caption{连续预留与 PagedAttention 分配 token slots 对比}
\label{tab:contiguous-vs-paged-slots}
\begin{tabular}{rrrrr}
\toprule
\textbf{请求} & \textbf{实际长度} & \textbf{连续预留} & \textbf{Paged block size 16} & \textbf{Paged 浪费} \
\midrule
1 & 130 & 4096 & 144 & 14 \
2 & 257 & 4096 & 272 & 15 \
3 & 519 & 4096 & 528 & 9 \
4 & 2050 & 4096 & 2064 & 14 \
5 & 4097 & 4096 & 4112 & 15 \
\bottomrule
\end{tabular}
\end{table}

表中第 5 个请求超过了 $4096$，说明连续预留若固定为 $4096$ 甚至无法容纳；若把 $S_{\max}$ 提高到更大，短请求浪费会更严重。Paged allocation 则可以随请求增长追加 blocks。

\subsection{吞吐提升}
\label{subsec:pagedattention-throughput-improvement}

PagedAttention 提升吞吐的核心路径可以写成：

[
\text{更高 KV 利用率}
\rightarrow
\text{同样显存容纳更多请求}
\rightarrow
\text{decode batch 更大}
\rightarrow
\text{GPU 利用率更高}
\rightarrow
\text{tokens/s 提升}.
]

Decode 阶段每个请求每次只生成一个 token。如果 batch 太小，GEMM 和 attention kernel 难以充分利用 GPU。通过降低 KV Cache 浪费，系统可以保留更多 running requests，从而提高 decode batch size。

但吞吐提升不是无限的。Decode batch 变大后，KV Cache 读取量也增加。最终系统可能受 HBM 带宽、attention kernel、sampling、CPU scheduler 或网络输出限制。PagedAttention 解决的是 KV Cache 分配和调度瓶颈，不会消除 decode attention 本身随上下文长度增长的读取成本。

可以把单 GPU serving 吞吐粗略看作：

[
\operatorname{Throughput}
\approx
\frac{
\operatorname{ActiveRequests}
}
{
T_{\mathrm{iteration}}
}.
]

PagedAttention 增加的是可同时活跃的请求数，但 $T_{\mathrm{iteration}}$ 也可能随着 batch 和上下文增长而上升。实际收益取决于请求长度分布、模型大小、GPU 显存、block size、kernel 实现和 scheduler 策略。

\subsection{长短请求混部}
\label{subsec:mixed-long-short-requests}

长短请求混部是在线 serving 的常态。PagedAttention 对混部尤其重要，因为它避免了短请求被长请求的最大长度预留拖累。

在连续分配中，如果 batch 内允许一个长请求达到 $S_{\max}=32768$，那么短请求也可能被迫预留接近这个上限。这样会严重降低短请求并发能力。

在 paged allocation 中，短请求只分配少量 blocks；长请求随着生成或 prefill 逐步占用更多 blocks。两类请求可以更自然地共存。

这对实际服务有几个影响：

\begin{itemize}
\item 短请求不会因为系统支持长上下文而付出巨大显存预留；
\item 长请求不会要求一开始就占用所有最大上下文 cache；
\item 请求结束后 blocks 可立即回收；
\item scheduler 可以根据剩余 blocks 做 admission control；
\item prefix cache 可让多个长请求共享共同前缀；
\item chunked prefill 可减少超长 prompt 对 decode 的阻塞。
\end{itemize}

\begin{figure}[H]
\centering
\begin{tikzpicture}[
font=\small,
cell/.style={draw, minimum width=0.35cm, minimum height=0.32cm},
long/.style={fill=green!12},
short/.style={fill=blue!12},
free/.style={fill=gray!16}
]

```
\node[font=\bfseries] at (0,3.0) {PagedAttention 支持长短请求混部};

\node[anchor=east] at (-4.8,1.55) {Short reqs};
\foreach \r/\y/\len in {1/1.75/3,2/1.25/2,3/0.75/4} {
  \foreach \i in {0,...,5} {
    \ifnum\i<\len
      \node[cell, short] at ({-4.2+0.38*\i},\y) {};
    \else
      \node[cell, free] at ({-4.2+0.38*\i},\y) {};
    \fi
  }
}

\node[anchor=east] at (0.5,1.25) {Long req};
\foreach \i in {0,...,17} {
  \node[cell, long] at ({1.0+0.38*\i},1.25) {};
}

\node[align=left, text width=10.4cm] at (0,-0.8) {
  短请求只占用少量 blocks；长请求占用更多 blocks。
  二者不需要按同一个最大长度预留，因此显存利用率更高。
};
```

\end{tikzpicture}
\caption{长短请求在 paged KV Cache 中自然共存}
\label{fig:mixed-long-short-requests-paged-kv}
\end{figure}

\subsection{trade-off 与适用边界}
\label{subsec:pagedattention-tradeoffs-boundaries}

PagedAttention 不是免费的。它用间接寻址和更复杂的内存管理换取更高显存利用率。主要 trade-off 包括：

\begin{itemize}
\item \textbf{地址翻译开销}：kernel 需要通过 block table 找到 physical block；
\item \textbf{metadata 开销}：需要维护 block table、refcount、request state；
\item \textbf{kernel 复杂度}：attention 读取不再是简单连续地址；
\item \textbf{block size 选择困难}：太大浪费多，太小 metadata 和寻址开销大；
\item \textbf{最后 block 仍有碎片}：内部碎片上界为 block size；
\item \textbf{prefill 长 prompt 仍然昂贵}：PagedAttention 不改变 $S^2$ prefill attention 成本；
\item \textbf{decode 长上下文仍然 memory-bound}：它不消除读取历史 KV 的线性成本；
\item \textbf{多 GPU 下更复杂}：TP/PP、prefix cache、block table 同步需要额外设计。
\end{itemize}

Block size 是一个典型权衡。设 block size 为 $B_{\mathrm{blk}}$。较大的 block size：

\begin{itemize}
\item block table 更小；
\item kernel 内 block 处理更规整；
\item metadata 更少；
\item 但最后 block 内部碎片更大。
\end{itemize}

较小的 block size：

\begin{itemize}
\item 内部碎片更小；
\item 按需增长更精细；
\item 但 block table 更大；
\item 地址翻译更多；
\item kernel 可能更难高效向量化。
\end{itemize}

因此 PagedAttention 适合动态在线 serving、请求长度差异大、高并发、KV Cache 显存紧张的场景。对于离线批处理、请求长度固定、显存充足或 kernel 极端追求连续访存的场景，简单连续布局可能也有竞争力。

\section{PagedAttention 的工程实现细节}
\label{sec:pagedattention-engineering-details}

为了让本章更贴近工程实践，本节进一步展开 PagedAttention 落地时的几个关键问题：block size、物理内存布局、slot mapping、prefill 写入、decode 追加、以及与采样策略的交互。

\subsection{block size 选择}
\label{subsec:block-size-choice}

Block size 是 PagedAttention 的核心超参数之一。它决定了每个 physical block 能存多少 token 的 K/V。设 block size 为 $B_{\mathrm{blk}}$，每个 token 的 KV Cache 大小为：

[
M_{\mathrm{token}}
==================

2
\cdot
L
\cdot
n_{kv}
\cdot
d_h
\cdot
s_{\mathrm{dtype}}.
]

如果每个 block 覆盖一个请求的 $B_{\mathrm{blk}}$ 个 token，则每个请求级 block 的 KV 大小为：

[
M_{\mathrm{block}}
==================

B_{\mathrm{blk}}
\cdot
M_{\mathrm{token}}.
]

实际系统中，物理 KV block 可能按 layer、head 或 worker 内部 layout 切分，因此实现细节会不同。但从逻辑上看，block size 越大，单次分配粒度越粗。

选择 block size 时需要同时考虑：

\begin{itemize}
\item 平均请求长度；
\item 最后 block 内部碎片；
\item block table metadata 大小；
\item attention kernel block 内连续读取效率；
\item GPU memory alignment；
\item prefix sharing 粒度；
\item copy-on-write 粒度；
\item 调度器分配和释放开销。
\end{itemize}

一个简单的期望碎片估算是：如果请求长度对 block size 的余数近似均匀分布，则每个请求最后 block 的平均浪费约为：

[
\mathbb{E}[W]
\approx
\frac{B_{\mathrm{blk}}-1}{2}.
]

因此 block size 从 $16$ 增加到 $32$，平均每请求浪费 token slots 约增加一倍。但如果 block size 太小，kernel 每次处理的连续 token 过少，metadata 和地址翻译开销上升。

\subsection{physical block layout}
\label{subsec:physical-block-layout}

Paged KV Cache 的物理 layout 需要服务 attention kernel。一个可能的 layout 是：

[
\mat{K}*{blocks}
\in
\mathbb{R}^{
N*{\mathrm{blocks}}
\times
n_{kv}
\times
B_{\mathrm{blk}}
\times
d_h
}.
]

V Cache 同理。也可以把 layer 维放在外面：

[
[L,N_{\mathrm{blocks}},n_{kv},B_{\mathrm{blk}},d_h].
]

或者为每层单独维护 block pool。不同 layout 的取舍包括：

\begin{itemize}
\item 每层独立 block pool 管理简单，但 metadata 多；
\item 所有层合并的 block 可以减少请求级 metadata，但 block 更大；
\item head 维连续有利于某些向量化加载；
\item block 内 token 维连续有利于顺序扫描历史位置；
\item head_dim 维连续通常有利于 vectorized load。
\end{itemize}

高性能 kernel 通常会对 layout 做非常具体的优化。本书不展开某个实现的所有细节，但读者应记住：PagedAttention 的设计不是只有 block table，\textbf{block table、memory layout 和 attention kernel 必须一起设计}。

\subsection{slot mapping}
\label{subsec:slot-mapping}

在每个 iteration 中，scheduler 选出一批要执行的 tokens。每个 token 的新 K/V 需要写入某个 physical block 的某个 offset。为了让 GPU kernel 知道写到哪里，系统通常会构造 slot mapping。

对一个 token，slot mapping 记录：

[
\text{token}
\rightarrow
(\text{physical block id},\ \text{offset}).
]

对于 decode，每个请求通常追加一个 token。设请求当前长度为 $S_i$，新 token 位置为 $S_i$。则：

[
b=\left\lfloor\frac{S_i}{B_{\mathrm{blk}}}\right\rfloor,
\quad
o=S_i\bmod B_{\mathrm{blk}}.
]

若 $o=0$，说明需要申请新 block。之后从 block table 得到 physical block id：

[
p=\operatorname{BlockTable}[i,b].
]

Slot mapping 为：

[
(p,o).
]

对于 prefill，一个请求一次写入多个 token。系统可以为这些 token 批量生成 slot mapping。若 prompt 跨多个 blocks，slot mapping 会跨越多个 physical blocks。

\begin{lstlisting}[language=Python,caption={为 decode token 构造 slot mapping 的简化伪代码},label={lst:decode-slot-mapping-pseudocode}]
def build_decode_slot_mapping(requests, block_manager, block_size):
slot_mapping = []

```
for req in requests:
    pos = req.current_length
    logical_block = pos // block_size
    offset = pos % block_size

    if offset == 0:
        block_manager.append_slot(req.request_id, pos)

    physical_block = req.block_table[logical_block]
    slot_mapping.append((physical_block, offset))

return slot_mapping
```

\end{lstlisting}

Slot mapping 是连接 scheduler、block manager 和 model runner 的关键桥梁。Scheduler 只知道请求要生成 token；Block Manager 知道请求的 blocks；Model Runner 需要具体写入地址。Slot mapping 把这三者连接起来。

\subsection{prefill 与 decode 的不同路径}
\label{subsec:prefill-and-decode-different-paths}

Prefill 和 decode 对 PagedAttention 的要求不同。

Prefill 中，一个请求可能一次写入很多 token。系统需要：

\begin{itemize}
\item 为 prompt 长度申请足够 blocks；
\item 批量写入 K/V；
\item 对 prompt 内部执行 causal attention；
\item 生成最后位置 logits；
\item 处理长 prompt 的 chunking。
\end{itemize}

Decode 中，每个请求每 iteration 通常只写入一个 token。系统需要：

\begin{itemize}
\item 检查当前 block 是否有空位；
\item 必要时申请新 block；
\item 读取完整历史 blocks；
\item 写入新 token K/V；
\item 执行 sampling；
\item 更新请求长度和 stop 状态。
\end{itemize}

因此，serving engine 可能为 prefill 和 decode 使用不同的 attention kernel 或不同调度策略。Prefill 更接近大矩阵计算，decode 更受 KV 读取和 batch 组织影响。PagedAttention 主要解决两者共同的 KV Cache 存储问题，但它在 decode 阶段的收益往往更直接，因为 decode 请求持续增长、加入和退出最频繁。

\section{本章小结}
\label{sec:chapter13-summary}

本章系统介绍了 PagedAttention 与 vLLM 风格 serving 架构。核心思想可以概括为：\textbf{把 KV Cache 从连续大块分配改造成分页式 block 管理，并让 attention kernel 通过 block table 读取逻辑连续、物理离散的 K/V。}

\begin{itemize}
\item \textbf{LLM serving 中的 KV Cache 是动态内存管理问题}：请求长度不同、到达时间不同、结束时间不同，cache 会不断增长和释放。
\item \textbf{连续 KV Cache 分配会产生 over-reservation}：如果每个请求都按最大上下文长度预留，短请求会浪费大量显存。
\item \textbf{Block-based allocation 将浪费限制在最后一个 block}：若 block size 为 $B_{\mathrm{blk}}$，每个请求最多浪费 $B_{\mathrm{blk}}-1$ 个 token slots。
\item \textbf{PagedAttention 类似操作系统分页}：请求逻辑 token 序列对应 logical blocks，物理 KV Cache 对应 physical blocks，block table 类似 page table。
\item \textbf{Block table 是核心元数据}：它记录每个请求的 logical block 到 physical block 的映射，使请求无需连续显存。
\item \textbf{Attention kernel 需要通过 block table 读取 KV}：逻辑位置 $s$ 被分解为 logical block id 和 offset，再映射到 physical block 地址。
\item \textbf{共享 prefix 可以通过共享 physical blocks 实现}：多个请求的 block table 指向相同 prefix blocks，并使用 refcount 管理生命周期。
\item \textbf{Copy-on-write 支持分叉场景}：beam search、parallel sampling、speculative decoding 等场景可以共享前缀，只在写入分叉时复制 block。
\item \textbf{vLLM 风格架构由 scheduler、block manager、worker、model runner 等组件组成}：scheduler 负责 iteration-level 调度，block manager 负责 KV blocks，worker 在 GPU 上执行模型。
\item \textbf{Continuous batching 与 PagedAttention 相互配合}：前者让请求动态加入和退出，后者让这种动态调度不需要移动连续 KV Cache。
\item \textbf{PagedAttention 的吞吐收益主要来自更高显存利用率和更大有效并发}，而不是改变 attention 的数学 FLOPs。
\item \textbf{PagedAttention 也有 trade-off}：block table 间接寻址、metadata 管理、kernel 复杂度和 block size 选择都需要工程权衡。
\item \textbf{Prefill 与 decode 的需求不同}：prefill 更偏大块计算和批量写入，decode 更偏动态追加、历史 KV 读取和调度效率。
\end{itemize}

下一章将进入 FlashAttention。PagedAttention 解决的是 KV Cache 的内存管理和动态调度问题；FlashAttention 解决的是 Attention 计算中的 IO 瓶颈。我们会从标准 Attention 需要物化 $S\times S$ attention matrix 的问题出发，推导 online softmax，理解为什么把 Q/K/V 分块搬入 SRAM 计算可以显著减少 HBM 读写，并最终形成高性能 fused attention kernel。
